% Blueprint: the Euler–Maclaurin summation formula with remainder, and its
% specialization to the Riemann zeta function.
%
% This file is a leanblueprint content fragment. It uses the standard
% leanblueprint macros: \lean{} records the intended Lean declaration name,
% \uses{} records dependency edges (by \label), and \leanok marks a statement
% (or proof) whose Lean counterpart is complete.

\chapter{Euler--Maclaurin summation with remainder}
\label{cha:euler_maclaurin}

\section{Prerequisites from mathlib}
\label{sec:em_prereqs}

All of the following are already in mathlib and are used as black boxes. They
are recorded here only to fix notation and to make the dependency edges of the
new results explicit.

\begin{definition}[Bernoulli numbers]
\label{def:bernoulli}
\lean{bernoulli}
\leanok
The Bernoulli numbers $B_n \in \mathbb{Q}$, with the convention $B_1 = +\tfrac12$
(mathlib's \texttt{bernoulli}; the alternative convention is \texttt{bernoulli'}).
Only even-index values $B_{2k}$ are used below, where the two conventions agree.
\end{definition}

\begin{definition}[Bernoulli polynomials as functions]
\label{def:bernoulliFun}
\lean{bernoulliFun}
\leanok
For $k \in \mathbb{N}$, $\operatorname{bernoulliFun} k : \mathbb{R} \to \mathbb{R}$
is the $k$-th Bernoulli polynomial $B_k(x)$ evaluated as a real function. mathlib
provides: $B_0(x) = 1$ (\texttt{bernoulliFun\_zero}), $B_1(x) = x - \tfrac12$
(\texttt{bernoulliFun\_one}), the derivative law
$\frac{d}{dx} B_k(x) = k\,B_{k-1}(x)$ (\texttt{hasDerivAt\_bernoulliFun}),
smoothness (\texttt{contDiff\_bernoulliFun}), the endpoint relation
$B_k(1) = B_k(0)$ for $k \neq 1$ (\texttt{bernoulliFun\_endpoints\_eq\_of\_ne\_one}),
$B_k(0) = B_k$ (\texttt{bernoulliFun\_eval\_zero}), and the vanishing mean
$\int_0^1 B_k = 0$ for $k \neq 0$ (\texttt{integral\_bernoulliFun\_eq\_zero}).
\end{definition}

\begin{definition}[Periodized Bernoulli functions]
\label{def:periodizedBernoulli}
\lean{periodizedBernoulli}
\leanok
For $k \in \mathbb{N}$, $\widetilde{B_k} := \operatorname{periodizedBernoulli} k$
is the $1$-periodic function on $\mathbb{R}$ (formally a function on
\texttt{UnitAddCircle}) that agrees with $B_k(x)$ on $[0,1)$. Equivalently
$\widetilde{B_k}(x) = B_k(\operatorname{Int.fract} x)$. It is the saw-tooth kernel
of the Euler--Maclaurin remainder. mathlib (\texttt{Mathlib.NumberTheory.ZetaValues})
provides its Fourier coefficients and basic regularity.
\end{definition}

\begin{theorem}[Integration by parts; mathlib]
\label{thm:ibp}
\lean{intervalIntegral.integral_mul_deriv_eq_deriv_mul}
\leanok
For $u, v$ with $u', v'$ interval-integrable on $[a,b]$,
$\int_a^b u\,v' = u(b)v(b) - u(a)v(a) - \int_a^b u'\,v$.
\end{theorem}

\begin{proof}
\leanok
In Mathlib as \texttt{intervalIntegral.integral\_mul\_deriv\_eq\_deriv\_mul}.
\end{proof}

\section{The base case: trapezoidal rule with remainder}
\label{sec:em_base}

\begin{lemma}[Unit-interval Euler--Maclaurin step]
\label{lem:em_step}
\lean{eulerMaclaurin_integral_unitInterval_eq}
\leanok
\uses{def:bernoulliFun, thm:ibp}
Let $f : \mathbb{R} \to \mathbb{R}$ be continuously differentiable on $[0,1]$.
Then
\[
  \int_0^1 f(x)\,dx
    \;=\; \frac{f(0)+f(1)}{2} \;-\; \int_0^1 \Bigl(x-\tfrac12\Bigr) f'(x)\,dx
    \;=\; \frac{f(0)+f(1)}{2} \;-\; \int_0^1 B_1(x)\,f'(x)\,dx .
\]
\end{lemma}

\begin{proof}
\leanok
\uses{thm:ibp, def:bernoulliFun}
Apply integration by parts (\Cref{thm:ibp}) with $u(x) = B_1(x) = x-\tfrac12$ and
$v = f$, so $u'(x) = 1$ and $v' = f'$:
\[
  \int_0^1 B_1(x)\,f'(x)\,dx
  = B_1(1)f(1) - B_1(0)f(0) - \int_0^1 1\cdot f(x)\,dx
  = \tfrac12 f(1) + \tfrac12 f(0) - \int_0^1 f(x)\,dx,
\]
using $B_1(1) = \tfrac12$ and $B_1(0) = -\tfrac12$. Rearranging gives the claim.
The hypothesis that $f$ is $C^1$ on $[0,1]$ supplies the interval-integrability of
$u'=1$ and $v'=f'$ required by \Cref{thm:ibp}.
\end{proof}

\section{The remainder term}
\label{sec:em_remainder}

\begin{definition}[Euler--Maclaurin remainder]
\label{def:em_remainder}
\lean{eulerMaclaurinRemainder}
\leanok
\uses{def:periodizedBernoulli, def:bernoulliFun}
For $p \ge 1$, a function $f$ that is $C^p$ on $[a,b]$, define
\[
  R_p(f; a, b) \;:=\; \frac{(-1)^{p+1}}{p!}\int_a^b \widetilde{B_p}(x)\, f^{(p)}(x)\,dx,
\]
where $f^{(p)}$ is the $p$-th iterated derivative and $\widetilde{B_p}$ is the
periodized Bernoulli function of \Cref{def:periodizedBernoulli}. To pin down the
integrand as a function of the real variable $x$ (note $\widetilde{B_p}$ is
formally a function on \texttt{UnitAddCircle}), we read $\widetilde{B_p}(x)$ as
$\operatorname{bernoulliFun} p\,(\operatorname{Int.fract} x)$, which agrees with
\texttt{periodizedBernoulli} under the quotient $\mathbb{R} \to$
\texttt{UnitAddCircle}.
\end{definition}

\section{Order raising on a single unit interval}
\label{sec:em_tower}

The order of the remainder is raised one unit interval at a time, because on $[n,n+1]$
the periodized kernel $\widetilde{B_k}$ coincides with the Bernoulli polynomial
$x \mapsto B_k(x-n)$, which is $C^\infty$. Every integration by parts then satisfies the
hypotheses of \Cref{thm:ibp}. The periodized function $\widetilde{B_k}$ itself is only
$C^0$ (with corners at the integers), so \Cref{thm:ibp} does not apply to it directly on
$[a,b]$.

The tower is assembled from two building blocks: the base case
(\Cref{lem:em_step} translated to $[n,n+1]$) and the order-raising
integration-by-parts recurrence.

\begin{lemma}[Per-interval base step]
\label{lem:em_step_translate}
\lean{eulerMaclaurin_integral_interval_eq}
\leanok
\uses{lem:em_step, def:bernoulliFun, thm:ibp}
Let $n \in \mathbb{R}$ and let $f$ be continuously differentiable on $[n,n+1]$. Then
\[
  \int_n^{n+1} f(x)\,dx
    \;=\; \frac{f(n)+f(n+1)}{2} \;-\; \int_n^{n+1} B_1(x-n)\,f'(x)\,dx .
\]
This is \Cref{lem:em_step} with the unit-translated Bernoulli kernel $B_1(x-n)$;
the proof is the same single integration by parts, now with $u(x)=B_1(x-n)$
(so $u'(x)=1$) and the boundary values $B_1(1)=\tfrac12$, $B_1(0)=-\tfrac12$.
\end{lemma}

\begin{proof}
\leanok
\uses{lem:em_step, def:bernoulliFun, thm:ibp}
The same single integration by parts as \Cref{lem:em_step}, with $u(x)=B_1(x-n)$
(so $u'(x)=1$) and boundary values $B_1(1)=\tfrac12$, $B_1(0)=-\tfrac12$.
\end{proof}

\begin{lemma}[Order-raising recurrence]
\label{lem:em_recurrence}
\lean{eulerMaclaurin_integral_bernoulliFun_smul_eq}
\leanok
\uses{def:bernoulliFun, def:bernoulli, thm:ibp}
Let $n \in \mathbb{R}$, let $q \ge 1$, and let $g$ be continuously differentiable on
$[n,n+1]$. Then, writing $\Delta g := g(n+1)-g(n)$,
\[
  \int_n^{n+1} B_q(x-n)\,g(x)\,dx
    \;=\; \frac{B_{q+1}}{q+1}\,\Delta g
        \;-\; \frac{1}{q+1}\int_n^{n+1} B_{q+1}(x-n)\,g'(x)\,dx .
\]
\end{lemma}

\begin{proof}
\leanok
\uses{def:bernoulliFun, def:bernoulli, thm:ibp}
Integrate by parts (\Cref{thm:ibp}) with $u(x)=B_{q+1}(x-n)/(q+1)$ and $v=g$, so
that $u'(x)=B_q(x-n)$ (the derivative law $B_{q+1}'=(q+1)B_q$ of
\Cref{def:bernoulliFun}, divided by $q+1$) and $v'=g'$. Both factors are smooth on
$[n,n+1]$, so all four hypotheses of \Cref{thm:ibp} hold. The boundary term is
$\tfrac1{q+1}\bigl(B_{q+1}(1)g(n+1)-B_{q+1}(0)g(n)\bigr)$, and since $q+1 \ge 2 \neq 1$
we have $B_{q+1}(1)=B_{q+1}(0)=B_{q+1}$
(\texttt{bernoulliFun\_endpoints\_eq\_of\_ne\_one} and
\texttt{bernoulliFun\_eval\_zero}). Solving the integration-by-parts identity for
$\int_n^{n+1}B_q(x-n)g(x)\,dx$ gives the claim.
\end{proof}

\begin{lemma}[Per-interval Euler--Maclaurin tower, raw form]
\label{lem:em_unit_tower_raw}
\lean{eulerMaclaurin_integral_interval_eq_tower}
\leanok
\uses{lem:em_step_translate, lem:em_recurrence, def:bernoulliFun, def:bernoulli, thm:ibp}
Let $n \in \mathbb{R}$ and $p \ge 1$. In the role of the iterated derivatives
$f^{(q)}$, take an explicit family of functions $F_0, F_1, \dots, F_p$ on
$[n,n+1]$ with $F_q' = F_{q+1}$ (i.e.\ $F_{q+1}$ is the derivative of $F_q$) for
$0 \le q < p$, each $F_q$ ($1 \le q \le p$) interval-integrable. Then, writing
$\Delta F_q := F_q(n+1) - F_q(n)$,
\[
  \int_n^{n+1} F_0(x)\,dx
  \;=\; \frac{F_0(n)+F_0(n+1)}{2}
      \;+\; \sum_{q=1}^{p-1} \frac{(-1)^{q}\,B_{q+1}}{(q+1)!}\,\Delta F_q
      \;+\; \frac{(-1)^{p}}{p!}\int_n^{n+1} B_p(x-n)\, F_p(x)\,dx .
\]
This is the per-interval tower before the even/odd collapse; specializing $F_q$ to
$f^{(q)}$ recovers the iterated-derivative form. The collapse to even indices
(\Cref{lem:em_unit_tower}) is a separate cosmetic step using the vanishing of
odd-index Bernoulli numbers.
\end{lemma}

\begin{proof}
\leanok
\uses{lem:em_step_translate, lem:em_recurrence, def:bernoulliFun, def:bernoulli}
Induction on $p \ge 1$. \emph{Base $p = 1$:} the sum $\sum_{q=1}^{0}$ is empty and
$\tfrac{(-1)^1}{1!} = -1$, so the claim is exactly the per-interval base step
(\Cref{lem:em_step_translate}) applied to $f = F_0$, $f' = F_1$. \emph{Step
$p \to p+1$:} apply the induction hypothesis (its derivative/integrability premises
are the restrictions of those for $p+1$) to rewrite $\int_n^{n+1}F_0$, then expand the
leftover order-$p$ integral $T_p := \int_n^{n+1} B_p(x-n)F_p$ by the order-raising
recurrence (\Cref{lem:em_recurrence}) with $g = F_p$, $g' = F_{p+1}$:
\[
  \frac{(-1)^p}{p!}T_p
  = \frac{(-1)^p}{p!}\Bigl(\frac{B_{p+1}}{p+1}\Delta F_p - \frac{1}{p+1}T_{p+1}\Bigr)
  = \frac{(-1)^p B_{p+1}}{(p+1)!}\Delta F_p + \frac{(-1)^{p+1}}{(p+1)!}T_{p+1},
\]
using $p!\,(p+1) = (p+1)!$ and $-(-1)^p = (-1)^{p+1}$. The first summand is exactly the
new $q = p$ term of $\sum_{q=1}^{p}$, and the second is the new leftover integral, which
is the claim for $p+1$.
\end{proof}

\begin{lemma}[Global Euler--Maclaurin tower, raw form]
\label{lem:em_global_raw}
\lean{eulerMaclaurin_integral_eq_tower_sum}
\leanok
\uses{lem:em_unit_tower_raw, lem:em_step_translate, lem:em_recurrence, def:bernoulliFun, def:bernoulli}
Let $a \in \mathbb{R}$, $N \in \mathbb{N}$, $p \ge 1$, and let $F_0, \dots, F_p$ be a
family on $[a, a+N]$ with $F_q' = F_{q+1}$ for $0 \le q < p$ and each $F_q$
($1 \le q \le p$) interval-integrable. Then
\[
  \int_a^{a+N} F_0(x)\,dx
  \;=\; \sum_{n=0}^{N-1} \frac{F_0(a+n)+F_0(a+n+1)}{2}
      \;+\; \sum_{q=1}^{p-1} \frac{(-1)^{q}\,B_{q+1}}{(q+1)!}\,\bigl(F_q(a+N) - F_q(a)\bigr)
      \;+\; \frac{(-1)^{p}}{p!}\sum_{n=0}^{N-1}\int_{a+n}^{a+n+1} B_p(x-(a+n))\, F_p(x)\,dx .
\]
This is the per-interval raw tower (\Cref{lem:em_unit_tower_raw}) summed over the $N$ unit
intervals: the integrals concatenate, the Bernoulli boundary differences telescope, and
the order-$p$ remainder constant factors out. It is the global tower before the even/odd
collapse and before identifying the remainder integrals with the periodized-Bernoulli
integral of \Cref{def:em_remainder}.
\end{lemma}

\begin{proof}
\leanok
\uses{lem:em_unit_tower_raw}
Apply \Cref{lem:em_unit_tower_raw} on each $[a+n, a+n+1]$ (its derivative and
integrability premises are the restrictions of the global ones, since
$[a+n,a+n+1]\subseteq[a,a+N]$ for $n<N$) and sum over $n = 0, \dots, N-1$. The left-hand
sides concatenate to $\int_a^{a+N} F_0$ (mathlib
\texttt{intervalIntegral.sum\_integral\_adjacent\_intervals}, using continuity of $F_0$).
On the right, the trapezoidal terms collect directly; the order-$q$ boundary differences
$\sum_{n}\bigl(F_q(a+n+1)-F_q(a+n)\bigr)$ telescope to $F_q(a+N)-F_q(a)$ (mathlib
\texttt{Finset.sum\_range\_sub}) after exchanging the $n$- and $q$-summations; and the
constant $\tfrac{(-1)^p}{p!}$ factors out of the remainder sum.
\end{proof}

\begin{lemma}[Global Euler--Maclaurin summation, raw form]
\label{lem:em_sum_raw}
\lean{eulerMaclaurin_sum_eq_tower_integral}
\leanok
\uses{lem:em_global_raw}
Under the hypotheses of \Cref{lem:em_global_raw} (a family $F_0, \dots, F_p$ on $[a,a+N]$
with $F_q' = F_{q+1}$ for $0 \le q < p$ and each $F_q$, $1 \le q \le p$, interval-integrable),
\[
  \sum_{n=0}^{N} F_0(a+n)
  \;=\; \int_a^{a+N} F_0(x)\,dx \;+\; \frac{F_0(a)+F_0(a+N)}{2}
      \;-\; \sum_{q=1}^{p-1} \frac{(-1)^{q}\,B_{q+1}}{(q+1)!}\,\bigl(F_q(a+N) - F_q(a)\bigr)
      \;-\; \frac{(-1)^{p}}{p!}\sum_{n=0}^{N-1}\int_{a+n}^{a+n+1} B_p(x-(a+n))\, F_p(x)\,dx .
\]
This is the value-sum counterpart of \Cref{lem:em_global_raw}: it is the global raw tower
rearranged so that the left-hand side is the sum of the $N+1$ sample values
$F_0(a), \dots, F_0(a+N)$. It is the raw form of the Euler--Maclaurin summation formula
(\Cref{thm:euler_maclaurin}), still before the even/odd Bernoulli collapse and before the
remainder integrals are identified with the periodized-Bernoulli integral
$R_p$ of \Cref{def:em_remainder}.
\end{lemma}

\begin{proof}
\leanok
\uses{lem:em_global_raw}
Rearrange \Cref{lem:em_global_raw} for the trapezoidal sum and apply the trapezoidal
reindexing
\[
  \sum_{n=0}^{N-1} \frac{F_0(a+n)+F_0(a+n+1)}{2}
    \;=\; \sum_{n=0}^{N} F_0(a+n) \;-\; \frac{F_0(a)+F_0(a+N)}{2},
\]
which is an immediate induction on $N$ (the half-endpoints telescope). Substituting and
solving for $\sum_{n=0}^{N} F_0(a+n)$ gives the claim.
\end{proof}

\section{The even/odd collapse}
\label{sec:em_collapse}

\begin{lemma}[Even/odd collapse of the Bernoulli boundary sum]
\label{lem:bernoulli_even_collapse}
\lean{eulerMaclaurin_sum_Ico_bernoulli_collapse}
\leanok
\uses{def:bernoulli}
For any coefficient family $c : \mathbb{N} \to \mathbb{R}$ and any $p \in \mathbb{N}$,
\[
  \sum_{q=1}^{p-1} \frac{(-1)^{q} B_{q+1}}{(q+1)!}\, c(q)
  \;=\; -\sum_{k=1}^{\lfloor p/2 \rfloor} \frac{B_{2k}}{(2k)!}\, c(2k-1).
\]
This is a purely combinatorial identity in $c$; applied with
$c(q) = F_q(a+N) - F_q(a)$ it converts the signed alternating Bernoulli boundary sum of
the raw towers (\Cref{lem:em_global_raw}, \Cref{lem:em_sum_raw}) into the classical
even-index form.
\end{lemma}

\begin{proof}
\leanok
\uses{def:bernoulli}
Two facts drive the collapse. First, for every $q \ge 1$ one has
$(-1)^q B_{q+1} = -B_{q+1}$: if $q$ is odd this is immediate, and if $q$ is even then
$q+1$ is odd and $\ge 3$, so $B_{q+1} = 0$ (the odd-index Bernoulli numbers vanish,
\texttt{bernoulli\_eq\_zero\_of\_odd}) and both sides are $0$. Hence every summand equals
$-\frac{B_{q+1}}{(q+1)!}c(q)$, and the terms with $q$ even vanish. Discarding them
(\texttt{Finset.sum\_filter\_of\_ne}) leaves the sum over odd $q$, which we reindex by
$q = 2k-1$, $k = (q+1)/2$ (a bijection of $\{q \in [1,p) : q \text{ odd}\}$ with
$\{1, \dots, \lfloor p/2\rfloor\}$, \texttt{Finset.sum\_bij'}); the surviving term at
$q = 2k-1$ is exactly $-\frac{B_{2k}}{(2k)!}c(2k-1)$.
\end{proof}

\begin{lemma}[Global Euler--Maclaurin summation, even-index form]
\label{lem:em_sum_even}
\lean{eulerMaclaurin_sum_eq_tower_integral_even}
\leanok
\uses{lem:em_sum_raw, lem:bernoulli_even_collapse}
Under the hypotheses of \Cref{lem:em_sum_raw},
\[
  \sum_{n=0}^{N} F_0(a+n)
  \;=\; \int_a^{a+N} F_0(x)\,dx \;+\; \frac{F_0(a)+F_0(a+N)}{2}
      \;+\; \sum_{k=1}^{\lfloor p/2\rfloor} \frac{B_{2k}}{(2k)!}\,\bigl(F_{2k-1}(a+N) - F_{2k-1}(a)\bigr)
      \;-\; \frac{(-1)^{p}}{p!}\sum_{n=0}^{N-1}\int_{a+n}^{a+n+1} B_p(x-(a+n))\, F_p(x)\,dx .
\]
This is \Cref{lem:em_sum_raw} after the even/odd collapse
(\Cref{lem:bernoulli_even_collapse}) of its Bernoulli boundary sum; note the sign of the
boundary sum flips to $+$. The remainder integrals are still the per-interval ones, before
identification with the periodized-Bernoulli integral.
\end{lemma}

\begin{proof}
\leanok
\uses{lem:em_sum_raw, lem:bernoulli_even_collapse}
Apply the collapse (\Cref{lem:bernoulli_even_collapse}) with
$c(q) = F_q(a+N) - F_q(a)$ to the Bernoulli boundary sum of \Cref{lem:em_sum_raw}; the
overall minus sign in front of that sum turns the collapse's $-$ into $+$.
\end{proof}

\begin{lemma}[Per-interval Euler--Maclaurin tower]
\label{lem:em_unit_tower}
\lean{eulerMaclaurin_integral_unitInterval_eq_tower}
\leanok
\uses{lem:em_unit_tower_raw, lem:bernoulli_even_collapse, def:bernoulliFun, def:bernoulli, thm:ibp}
Let $n \in \mathbb{Z}$, $p \ge 1$, and let $f$ be $C^p$ on $[n, n+1]$. Writing
$\Delta g := g(n+1) - g(n)$,
\[
  \int_n^{n+1} f(x)\,dx
  \;=\; \frac{f(n)+f(n+1)}{2}
      \;-\; \sum_{k=1}^{\lfloor p/2\rfloor} \frac{B_{2k}}{(2k)!}\,\Delta f^{(2k-1)}
      \;+\; \frac{(-1)^{p}}{p!}\int_n^{n+1} B_p(x-n)\, f^{(p)}(x)\,dx .
\]
\end{lemma}

\begin{proof}
\leanok
\uses{lem:em_unit_tower_raw, lem:bernoulli_even_collapse, def:bernoulliFun, def:bernoulli}
The Lean proof reorganises the classical unrolling below: it specializes the raw
per-interval tower (\Cref{lem:em_unit_tower_raw}) to the iterated derivatives
$F_q = f^{(q)}$ and applies the even/odd collapse (\Cref{lem:bernoulli_even_collapse})
to its Bernoulli boundary sum.

Let $T_q := \int_n^{n+1} B_q(x-n)\,f^{(q)}(x)\,dx$ for $1 \le q \le p$, where
$x \mapsto B_q(x-n)$ is a smooth polynomial on $[n,n+1]$.

\emph{Base $p=1$.} \Cref{lem:em_step}, after the change of variables $x \mapsto x-n$
(a unit right-translation; mathlib \texttt{intervalIntegral.integral\_comp\_sub\_right})
and using $\widetilde{B_1}(x)=B_1(x-n)$ on $[n,n+1)$, reads
$\int_n^{n+1} f = \tfrac{f(n)+f(n+1)}{2} - T_1$. Since $\lfloor 1/2\rfloor = 0$ the
Bernoulli sum is empty and $\tfrac{(-1)^1}{1!}T_1 = -T_1$, so the claim holds.

\emph{Recurrence for $T_q$.} For $1 \le q$, integrate $T_q$ by parts
(\Cref{thm:ibp}) with $u(x) = B_{q+1}(x-n)/(q+1)$ and $v = f^{(q)}$, using
$\frac{d}{dx}\frac{B_{q+1}(x-n)}{q+1} = B_q(x-n)$ (the derivative law of
\Cref{def:bernoulliFun}). All four hypotheses of \Cref{thm:ibp} hold because both
factors are smooth on $[n,n+1]$. The boundary term is
$\frac{1}{q+1}\bigl(B_{q+1}(1) f^{(q)}(n+1) - B_{q+1}(0) f^{(q)}(n)\bigr)$, and since
$q+1 \ge 2 \neq 1$ we have $B_{q+1}(1) = B_{q+1}(0) = B_{q+1}$
(\texttt{bernoulliFun\_endpoints\_eq\_of\_ne\_one} and \texttt{bernoulliFun\_eval\_zero}).
Hence
\[
  T_q \;=\; \frac{B_{q+1}}{q+1}\,\Delta f^{(q)} \;-\; \frac{1}{q+1}\,T_{q+1}.
\]

\emph{Unrolling.} Substituting $T_1$ from the recurrence into the base identity and
iterating up to $q = p-1$, an immediate induction shows the coefficient of $T_q$ in
$\int_n^{n+1} f$ is $\frac{(-1)^q}{q!}$, and the accumulated boundary terms are
$-\sum_{q=1}^{p-1}\frac{(-1)^{q-1}}{q!}\cdot\frac{B_{q+1}}{q+1}\Delta f^{(q)}
 = -\sum_{q=1}^{p-1}\frac{(-1)^{q-1}B_{q+1}}{(q+1)!}\Delta f^{(q)}$.
Now $(-1)^{q-1}B_{q+1} = B_{q+1}$ for every $q \ge 1$: if $q+1$ is even then
$(-1)^{q-1} = +1$; if $q+1$ is odd and $\ge 3$ then $B_{q+1} = 0$
(odd-index Bernoulli numbers vanish), so the equality holds trivially. Therefore
only the terms with $q+1 = 2k$ survive, contributing exactly
$-\sum_{k=1}^{\lfloor p/2\rfloor}\frac{B_{2k}}{(2k)!}\Delta f^{(2k-1)}$, and the
leftover term is $\frac{(-1)^p}{p!}T_p$, as claimed. \qed
\end{proof}

\section{The Euler--Maclaurin summation formula}
\label{sec:em_main}

\begin{theorem}[Euler--Maclaurin summation formula with remainder]
\label{thm:euler_maclaurin}
\lean{eulerMaclaurin_sum_eq}
\leanok
\uses{def:bernoulli, def:em_remainder, lem:em_unit_tower, lem:em_sum_even, def:periodizedBernoulli}
Let $a \le b$ be integers, let $p \ge 1$, and let $E$ be a real Banach space, and let
$f : \mathbb{R} \to E$ be $C^p$ on $[a,b]$. Then
\[
  \sum_{n=a}^{b} f(n)
  \;=\; \int_a^b f(x)\,dx
      \;+\; \frac{f(a)+f(b)}{2}
      \;+\; \sum_{k=1}^{\lfloor p/2\rfloor} \frac{B_{2k}}{(2k)!}\Bigl(f^{(2k-1)}(b) - f^{(2k-1)}(a)\Bigr)
      \;+\; R_p(f; a, b).
\]
\end{theorem}

\begin{proof}
\leanok
\uses{lem:em_sum_even, def:em_remainder, def:bernoulliFun, def:periodizedBernoulli}
In Lean the global summation is carried out once at the family level
(\Cref{lem:em_sum_even}), specialized here to $F_q = f^{(q)}$; the per-interval
identity below (\Cref{lem:em_unit_tower}) is its mathematical content. It then only
remains to identify the per-interval remainder integrals with $R_p$: with $a$ an
integer, on $[a+n,a+n+1)$ the kernel $B_p(x-(a+n))$ equals $B_p(\operatorname{Int.fract} x)$
almost everywhere, so the unit integrals concatenate to
$\int_a^{a+N} \widetilde{B_p}(x) f^{(p)}(x)\,dx$.

Sum the per-interval identity of \Cref{lem:em_unit_tower} over the integer points
$n = a, a+1, \dots, b-1$ (a $C^p$ function on $[a,b]$ is $C^p$ on each subinterval):
\[
  \int_a^b f
  = \sum_{n=a}^{b-1}\frac{f(n)+f(n+1)}{2}
    - \sum_{k=1}^{\lfloor p/2\rfloor}\frac{B_{2k}}{(2k)!}\sum_{n=a}^{b-1}\Delta f^{(2k-1)}
    + \frac{(-1)^p}{p!}\sum_{n=a}^{b-1}\int_n^{n+1} B_p(x-n)\,f^{(p)}(x)\,dx .
\]
Three telescoping/concatenation facts finish the proof.
\begin{enumerate}
\item \emph{Trapezoidal endpoints:}
  $\sum_{n=a}^{b-1}\frac{f(n)+f(n+1)}{2} = \sum_{n=a}^{b} f(n) - \frac{f(a)+f(b)}{2}$.
\item \emph{Bernoulli differences telescope:}
  $\sum_{n=a}^{b-1}\Delta f^{(2k-1)} = \sum_{n=a}^{b-1}\bigl(f^{(2k-1)}(n+1)-f^{(2k-1)}(n)\bigr)
   = f^{(2k-1)}(b)-f^{(2k-1)}(a)$.
\item \emph{Remainder integrals concatenate:} on $[n,n+1)$ one has
  $B_p(x-n) = \widetilde{B_p}(x)$ (\Cref{def:periodizedBernoulli}), so the unit
  integrals join to $\int_a^b \widetilde{B_p}(x) f^{(p)}(x)\,dx$, whence
  $\frac{(-1)^p}{p!}\sum_n \int_n^{n+1}B_p(x-n)f^{(p)}
  = -\frac{(-1)^{p+1}}{p!}\int_a^b \widetilde{B_p}f^{(p)} = -R_p(f;a,b)$
  (\Cref{def:em_remainder}).
\end{enumerate}
Substituting (1)--(3) gives
$\int_a^b f = \sum_{n=a}^b f(n) - \frac{f(a)+f(b)}{2}
 - \sum_{k}\frac{B_{2k}}{(2k)!}\bigl(f^{(2k-1)}(b)-f^{(2k-1)}(a)\bigr) - R_p(f;a,b)$,
and rearranging for $\sum_{n=a}^b f(n)$ is exactly the assertion. \qed
\end{proof}

\begin{remark}
\label{rem:em_side_conditions}
\uses{thm:euler_maclaurin}
The side conditions needed to discharge \Cref{thm:euler_maclaurin} in Lean are:
(i) interval-integrability of each $\widetilde{B_p}\,f^{(p)}$, which follows from
continuity of $\widetilde{B_p}$ for $p \ge 2$ and boundedness for $p=1$ together
with $C^p$ regularity of $f$; (ii) the iterated-derivative bookkeeping relating
$f^{(p)}$ to \texttt{iteratedDeriv}. These are routine but voluminous.
\end{remark}

\section{Complex-valued Euler--Maclaurin}
\label{sec:em_complex}

The zeta application takes $f(x) = x^{-s}$, a $\mathbb{C}$-valued function, whereas the
classical statement is for $f : \mathbb{R} \to \mathbb{R}$. No transfer argument is needed.
The periodized Bernoulli kernel $\widetilde{B_p}$ is a \emph{real} scalar, acting on $f^{(p)}$
by scalar multiplication, and the Bernoulli coefficients $B_{2k}/(2k)!$ are real. So both
\Cref{def:em_remainder} and \Cref{thm:euler_maclaurin} are stated once for $f : \mathbb{R} \to
E$ with $E$ any real Banach space, the integral being the Bochner integral. The complex case is
the instance $E = \mathbb{C}$.

\begin{lemma}[Complex-valued Euler--Maclaurin summation]
\label{lem:em_sum_complex}
\lean{Complex.sum_eq}
\leanok
\uses{thm:euler_maclaurin, def:em_remainder}
Let $a$ be an integer, $p \ge 1$, $N \in \mathbb{N}$ and $f : \mathbb{R} \to \mathbb{C}$ be such
that the iterated derivatives satisfy the derivative chain $(f^{(q)})' = f^{(q+1)}$ on
$[a, a+N]$ for $q < p$, and each $f^{(q)}$ ($1 \le q \le p$) is interval-integrable on
$[a, a+N]$. Then
\[
  \sum_{n=0}^{N} f(a+n)
  \;=\; \int_a^{a+N} f(x)\,dx \;+\; \frac{f(a)+f(a+N)}{2}
      \;+\; \sum_{k=1}^{\lfloor p/2\rfloor} \frac{B_{2k}}{(2k)!}\bigl(f^{(2k-1)}(a+N) - f^{(2k-1)}(a)\bigr)
      \;+\; R_p(f; a, a+N),
\]
the verbatim $\mathbb{C}$-valued analogue of \Cref{thm:euler_maclaurin}, with the real scalar
action read as complex multiplication.
\end{lemma}

\begin{proof}
\leanok
\uses{thm:euler_maclaurin}
Instantiate \Cref{thm:euler_maclaurin} at $E = \mathbb{C}$ and rewrite the real scalar action
as complex multiplication (\texttt{Complex.real\_smul}). The remainder is put into its
multiplicative form $R_p(f;a,b) = \frac{(-1)^{p+1}}{p!}\int_a^b \widetilde{B_p}(x) f^{(p)}(x)\,dx$
by \texttt{Complex.eulerMaclaurinRemainder\_eq}.
\end{proof}

\section{The power function \texorpdfstring{$x^{-s}$}{x^{-s}} and its derivatives}
\label{sec:em_power}

Throughout, $f_s : \mathbb{R} \to \mathbb{C}$ denotes $f_s(x) = (x:\mathbb{C})^{-s}$, and
$(s)_j := \prod_{i=0}^{j-1}(s+i)$ is the rising factorial (mathlib \texttt{ascPochhammer},
with $(s)_j = (\texttt{ascPochhammer}\ \mathbb{C}\ j).\mathrm{eval}\,s$; in particular
$(s)_0 = 1$ and $(s)_{j+1} = (s)_j\,(s+j)$).

\begin{lemma}[Iterated derivatives of $x^{-s}$]
\label{lem:iteratedDeriv_cpow}
\lean{Complex.iteratedDeriv_cpow_ofReal_neg}
\leanok
For $s \in \mathbb{C}$, $j \in \mathbb{N}$, and $x > 0$,
\[
  f_s^{(j)}(x) \;=\; (-1)^j\,(s)_j\,(x:\mathbb{C})^{-s-j}.
\]
The induction below in fact establishes the stronger \texttt{HasDerivAt} form
$\mathrm{HasDerivAt}\,(f_s^{(j)})\,\bigl((-1)^{j+1}(s)_{j+1}(x)^{-s-(j+1)}\bigr)\,x$ on
$x > 0$ (formalized as the companion lemma
\texttt{Complex.hasDerivAt\_iteratedDeriv\_cpow\_ofReal\_neg}, stated as
$\mathrm{HasDerivAt}\,(f_s^{(j)})\,(f_s^{(j+1)}(x))\,x$), which is what discharges the
derivative-chain hypothesis of \Cref{lem:em_sum_complex} in \Cref{lem:sum_cpow}.
\end{lemma}

\begin{proof}
\leanok
Induction on $j$. The single-step rule is, for $x > 0$ and any exponent
$c \in \mathbb{C}$, $\frac{d}{dx}(x:\mathbb{C})^{c} = c\,(x:\mathbb{C})^{c-1}$ as a
\texttt{HasDerivAt} statement (mathlib \texttt{hasDerivAt\_ofReal\_cpow\_const} for
$c \neq 0$; for $c = 0$ the function is the constant $1$ and both sides vanish). The base
$j = 0$ is $f_s(x) = (x)^{-s}$ with $(s)_0 = 1$. For the step, the open half-line
$x > 0$ lets us differentiate $x \mapsto (-1)^j(s)_j(x)^{-s-j}$ pointwise: the rule gives
\[
  (-1)^j(s)_j\,(-s-j)\,(x)^{-s-j-1}
  = (-1)^{j+1}(s)_j(s+j)\,(x)^{-s-(j+1)}
  = (-1)^{j+1}(s)_{j+1}\,(x)^{-s-(j+1)},
\]
using $-(-s-j) = s+j$ and $(s)_j(s+j) = (s)_{j+1}$. Combined with
$f_s^{(j+1)} = (f_s^{(j)})'$ (\texttt{iteratedDeriv\_succ}) this is the claim for $j+1$.
\end{proof}

\section{The finite Euler--Maclaurin expansion of the zeta partial sums}
\label{sec:em_zeta_finite}

\begin{lemma}[Finite Euler--Maclaurin expansion of $\sum n^{-s}$]
\label{lem:sum_cpow}
\lean{Complex.sum_cpow_eq}
\leanok
\uses{lem:em_sum_complex, lem:iteratedDeriv_cpow, def:em_remainder}
Let $s \in \mathbb{C}$ with $s \neq 1$, $M \in \mathbb{N}$, and integers $1 \le N \le R$.
Then, with $f_s(x) = (x:\mathbb{C})^{-s}$ and order $p = 2M+1$,
\[
  \sum_{n=N}^{R} n^{-s}
  = \frac{R^{1-s} - N^{1-s}}{1-s} + \frac{N^{-s}+R^{-s}}{2}
    - \sum_{k=1}^{M}\frac{B_{2k}}{(2k)!}(s)_{2k-1}\bigl(R^{-s-(2k-1)} - N^{-s-(2k-1)}\bigr)
    + R_{2M+1}(f_s; N, R).
\]
\end{lemma}

\begin{proof}
\leanok
\uses{lem:em_sum_complex, lem:iteratedDeriv_cpow}
Apply \Cref{lem:em_sum_complex} with $a = N$, the number of unit intervals $R - N$, order
$p = 2M+1$ (so $\lfloor p/2\rfloor = M$), and $f = f_s$. On $[N,R]$ with $N \ge 1$ the base
is positive, so \Cref{lem:iteratedDeriv_cpow} supplies the derivative chain
$(f_s^{(q)})' = f_s^{(q+1)}$, and each $f_s^{(q)}$ is continuous on $[N,R] \subset (0,\infty)$,
hence interval-integrable. Three substitutions finish.
\begin{enumerate}
\item \emph{Integral:} $\int_N^R (x)^{-s}\,dx = \frac{R^{1-s}-N^{1-s}}{1-s}$ for $s \neq 1$,
  since $\frac{d}{dx}\frac{x^{1-s}}{1-s} = (x)^{-s}$ on $x > 0$ (FTC, from the cpow
  derivative rule).
\item \emph{Bernoulli boundary terms:} by \Cref{lem:iteratedDeriv_cpow},
  $f_s^{(2k-1)}(x) = -(s)_{2k-1}(x)^{-s-(2k-1)}$ (as $(-1)^{2k-1} = -1$), so the $k$-th
  boundary term $\frac{B_{2k}}{(2k)!}\bigl(f_s^{(2k-1)}(R)-f_s^{(2k-1)}(N)\bigr)$ equals
  $-\frac{B_{2k}}{(2k)!}(s)_{2k-1}\bigl(R^{-s-(2k-1)}-N^{-s-(2k-1)}\bigr)$.
\item \emph{Endpoints:} $\frac{f_s(N)+f_s(R)}{2} = \frac{N^{-s}+R^{-s}}{2}$, and the
  reindexed left-hand side $\sum_{n=0}^{R-N} f_s(N+n) = \sum_{n=N}^{R} n^{-s}$.
\end{enumerate}
\end{proof}

\section{Specialization to the Riemann zeta function}
\label{sec:em_zeta}

Passing $R \to \infty$ in \Cref{lem:sum_cpow} and identifying the limit of the partial
sums with $\zeta(s)$ yields the Euler--Maclaurin expansion of $\zeta$. The limit is taken
first in the region of absolute convergence $\operatorname{Re} s > 1$
(\Cref{lem:zeta_em_conv}), then propagated to
$\operatorname{Re} s > -(2M+1)$, $s \neq 1$, by analytic continuation
(\Cref{lem:zeta_em_analytic} and \Cref{thm:zeta_em}). Write
\[
  R_M(s,N) \;:=\; R_{2M+1}(f_s; N, \infty)
  \;=\; \frac{1}{(2M+1)!}\int_N^{\infty} \widetilde{B}_{2M+1}(x)\,f_s^{(2M+1)}(x)\,dx
\]
for the improper order-$(2M+1)$ remainder (note $(-1)^{2M+2} = 1$). This integral converges
absolutely only for $\operatorname{Re} s > -2M$; on the remaining strip
$-(2M+1) < \operatorname{Re} s \le -2M$ the symbol $R_M(s,N)$ denotes instead the
analytically continued value supplied by the order-$(2M+2)$ representation of
\Cref{lem:zeta_em_analytic} (the two agree wherever both make sense).

\begin{definition}[Improper Euler--Maclaurin remainder]
\label{def:em_remainder_Ioi}
\lean{Complex.eulerMaclaurinRemainderIoi}
\leanok
\uses{def:periodizedBernoulli, def:em_remainder}
For $p \ge 1$ and $f$ that is $C^p$ on $[a,\infty)$, the improper remainder is
\[
  R_p^{\infty}(f; a) \;:=\; \frac{(-1)^{p+1}}{p!}\int_a^{\infty} \widetilde{B_p}(x)\, f^{(p)}(x)\,dx,
\]
the $b \to \infty$ limit of \Cref{def:em_remainder}, with $\widetilde{B_p}(x)$ read as
$\operatorname{bernoulliFun} p\,(\operatorname{Int.fract} x)$ as before. In the zeta
application $R_M(s,N) = R_{2M+1}^{\infty}(f_s; N)$.
\end{definition}

\begin{lemma}[Absolute convergence of the zeta remainder integral]
\label{lem:integrable_rem_zeta}
\lean{Complex.integrableOn_Ioi_bernoulliFun_fract_iteratedDeriv_cpow}
\leanok
\uses{def:periodizedBernoulli, lem:iteratedDeriv_cpow}
Let $N \ge 1$ and $\operatorname{Re} s + p > 1$. Then the integrand
$x \mapsto \widetilde{B_p}(x)\,f_s^{(p)}(x)$ of the order-$p$ remainder is (Lebesgue)
integrable on $(N, \infty)$. Indeed the periodized Bernoulli kernel is bounded, say
$|\widetilde{B_p}| \le C$ (it is $\operatorname{bernoulliFun} p$ composed with
$\operatorname{Int.fract}$, and $\operatorname{bernoulliFun} p$ is continuous on the compact
$[0,1]$), while $\|f_s^{(p)}(x)\| = |(s)_p|\,x^{-\operatorname{Re} s - p}$ by
\Cref{lem:iteratedDeriv_cpow}; hence $\|\widetilde{B_p}(x)\,f_s^{(p)}(x)\| \le
C\,|(s)_p|\,x^{-\operatorname{Re} s - p}$, and $x^{-\operatorname{Re} s - p}$ is integrable on
$(N,\infty)$ because $-\operatorname{Re} s - p < -1$. In particular $R_M(s,N)$ converges
absolutely for $\operatorname{Re} s > -2M$ (taking $p = 2M+1$), a fortiori for
$\operatorname{Re} s > 1$.
\end{lemma}

\begin{proof}
\leanok
\uses{lem:iteratedDeriv_cpow}
Dominated convergence: the integrand is a.e.-strongly-measurable on $(N,\infty)$ (the kernel
is measurable as $\operatorname{bernoulliFun} p \circ \operatorname{Int.fract}$, and
$x \mapsto x^{-s-p}$ is continuous on $(0,\infty) \supseteq (N,\infty)$), and is dominated by
the integrable $x \mapsto C\,|(s)_p|\,x^{-\operatorname{Re} s - p}$ as displayed
(\texttt{integrableOn\_Ioi\_rpow\_of\_lt} for the dominating function, using $N > 0$ and
$-\operatorname{Re} s - p < -1$). The pointwise bound uses
$\|(x:\mathbb{C})^{-s-p}\| = x^{-\operatorname{Re} s - p}$ for $x > 0$
(\texttt{Complex.norm\_cpow\_eq\_rpow\_re\_of\_pos}).
\end{proof}

\begin{lemma}[Remainder converges and the partial sums tend to $\zeta$, $\operatorname{Re} s > 1$]
\label{lem:zeta_em_conv}
\lean{riemannZeta_eq_of_one_lt_re}
\leanok
\uses{lem:sum_cpow, def:em_remainder_Ioi, lem:integrable_rem_zeta}
For $\operatorname{Re} s > 1$ and integer $N \ge 1$, the improper integral $R_M(s,N)$
converges and
\[
  \zeta(s) = \sum_{n=1}^{N-1} n^{-s} + \frac{N^{1-s}}{s-1} + \tfrac12 N^{-s}
    + \sum_{k=1}^{M} \frac{B_{2k}}{(2k)!}\,(s)_{2k-1}\,N^{-s-(2k-1)} + R_M(s,N).
\]
\end{lemma}

\begin{proof}
\leanok
\uses{lem:sum_cpow, def:em_remainder_Ioi, lem:integrable_rem_zeta}
Take $R \to \infty$ through the integers $R \ge N$ in \Cref{lem:sum_cpow}. The left-hand
partial sums $\sum_{n=N}^{R} n^{-s}$ converge to $\sum_{n=N}^{\infty} n^{-s} =
\zeta(s) - \sum_{n=1}^{N-1} n^{-s}$: the series $\sum_n n^{-s}$ is summable for
$\operatorname{Re} s > 1$ (\texttt{Complex.summable\_one\_div\_nat\_cpow}), its total sum is
$\zeta(s)$ (\texttt{zeta\_eq\_tsum\_one\_div\_nat\_add\_one\_cpow}, since the $n=0$ term
$0^{-s}$ vanishes), and $\sum_{n \in [N,R]} = \sum_{n<R+1} - \sum_{n<N}$ tends to the tail.
On the right-hand side, every power of $R$ has an exponent of negative real part
($\operatorname{Re}(1-s) < 0$, $\operatorname{Re}(-s) < 0$,
$\operatorname{Re}(-s-(2k-1)) < 0$), so $\|R^{c}\| = R^{\operatorname{Re} c} \to 0$
(\texttt{Complex.norm\_cpow\_eq\_rpow\_re\_of\_pos} and \texttt{tendsto\_rpow\_neg\_atTop});
hence $\frac{R^{1-s}-N^{1-s}}{1-s} \to \frac{-N^{1-s}}{1-s} = \frac{N^{1-s}}{s-1}$, the
endpoint term $\to \tfrac12 N^{-s}$, and the Bernoulli sum $\to +\sum_k \frac{B_{2k}}{(2k)!}
(s)_{2k-1} N^{-s-(2k-1)}$ (the $R$-half of each difference vanishes, flipping the sign). The
finite remainder $R_{2M+1}(f_s; N, R)$ converges to $R_M(s,N)$ because its integrand is
integrable on $(N,\infty)$ (\Cref{lem:integrable_rem_zeta}), so
$\int_N^R \to \int_N^{\infty}$ (\texttt{intervalIntegral\_tendsto\_integral\_Ioi}). Equating
the two limits (uniqueness of limits) and solving for $\zeta(s)$ gives the claim.
\end{proof}

The remainder $R_M(s,N) = R^{\infty}_{2M+1}(f_s; N)$ converges absolutely only for
$\operatorname{Re} s > -2M$ (\Cref{lem:integrable_rem_zeta} with $p = 2M+1$). To reach the
full strip $\operatorname{Re} s > -(2M+1)$ we replace the order-$(2M+1)$ integral by an
\emph{order-$(2M+2)$} representation obtained from one integration by parts; that integral
converges on the larger strip and is analytic there, which is what the identity theorem
needs. We first record the integration-by-parts identity, then the analyticity, then the
identity-theorem continuation.

\begin{lemma}[One integration by parts: order $2M+1 \to 2M+2$]
\label{lem:rem_ibp}
\lean{Complex.eulerMaclaurinRemainderIoi_succ_eq}
\leanok
\uses{def:em_remainder_Ioi, lem:integrable_rem_zeta, lem:iteratedDeriv_cpow, def:bernoulli, def:bernoulliFun, lem:em_recurrence, lem:em_global_raw}
Let $N \ge 1$ and $\operatorname{Re} s > -2M$ (so the order-$(2M+1)$ remainder integral
converges, \Cref{lem:integrable_rem_zeta}). Writing
\[
  t_{M+1}(s,N) := \frac{B_{2M+2}}{(2M+2)!}\,(s)_{2M+1}\,N^{-s-(2M+1)}
\]
for the first omitted even Bernoulli term, one has
\[
  R^{\infty}_{2M+1}(f_s; N) \;=\; t_{M+1}(s,N) \;+\; R^{\infty}_{2M+2}(f_s; N),
\]
where (recall $(-1)^{2M+2}=1$ and $(-1)^{2M+3}=-1$, so the two improper remainders are)
\[
  R^{\infty}_{2M+1}(f_s; N) = \frac{1}{(2M+1)!}\int_N^{\infty}\widetilde{B}_{2M+1}\,f_s^{(2M+1)},
  \qquad
  R^{\infty}_{2M+2}(f_s; N) = -\frac{1}{(2M+2)!}\int_N^{\infty}\widetilde{B}_{2M+2}\,f_s^{(2M+2)} .
\]
\end{lemma}

\begin{proof}
\leanok
\uses{def:em_remainder_Ioi, lem:integrable_rem_zeta, lem:iteratedDeriv_cpow, def:bernoulli, def:bernoulliFun, lem:em_recurrence, lem:em_global_raw}
Both integrals converge absolutely (\Cref{lem:integrable_rem_zeta}: the order-$(2M+1)$
integrand for $\operatorname{Re} s + (2M+1) > 1$, i.e.\ $\operatorname{Re} s > -2M$; the
order-$(2M+2)$ integrand for $\operatorname{Re} s + (2M+2) > 1$, a fortiori). It therefore
suffices to prove the displayed identity for the truncated integral on $[N,R]$ and pass
$R \to \infty$ (\texttt{intervalIntegral\_tendsto\_integral\_Ioi}, as in
\Cref{lem:zeta_em_conv}), the $R$-boundary contribution vanishing as below.

\emph{Integration by parts on $[N,R]$.} The kernel $\widetilde{B}_{2M+2}$ is $C^1$ on each
unit subinterval with $\frac{d}{dx}\widetilde{B}_{2M+2}(x) = (2M+2)\,\widetilde{B}_{2M+1}(x)$
away from the integers (the derivative law $B_{k}' = k B_{k-1}$ of \Cref{def:bernoulliFun}
applied to $\widetilde{B}_{k}(x) = B_k(\operatorname{Int.fract} x)$, whose jump-free
continuity for index $2M+2 \ge 2 \ne 1$ comes from
\texttt{bernoulliFun\_endpoints\_eq\_of\_ne\_one}); it is moreover globally continuous. Thus
on $[N,R]$, integrating $\int_N^R \widetilde{B}_{2M+1}\,f_s^{(2M+1)}$ by parts with
$u = \widetilde{B}_{2M+2}/(2M+2)$ (so $u' = \widetilde{B}_{2M+1}$) and $v = f_s^{(2M+1)}$
(so $v' = f_s^{(2M+2)}$, \Cref{lem:iteratedDeriv_cpow}) gives
\[
  \int_N^R \widetilde{B}_{2M+1}\,f_s^{(2M+1)}
  = \Bigl[\tfrac{\widetilde{B}_{2M+2}}{2M+2}\,f_s^{(2M+1)}\Bigr]_N^R
    - \frac{1}{2M+2}\int_N^R \widetilde{B}_{2M+2}\,f_s^{(2M+2)} .
\]
This is the per-unit-interval recurrence \Cref{lem:em_recurrence} summed over
$[N,R]$ exactly as in \Cref{lem:em_global_raw} (the interior boundary contributions
telescope and the periodized kernel is read on $[n,n+1)$ as $B_q(x-n)$), so no new
integration-by-parts hypothesis beyond those already discharged there is needed.

\emph{Boundary terms.} At an integer $x$, $\operatorname{Int.fract} x = 0$, so
$\widetilde{B}_{2M+2}(x) = B_{2M+2}(0) = B_{2M+2}$ (\texttt{bernoulliFun\_eval\_zero}). At
$x = N$ the boundary value is $\tfrac{B_{2M+2}}{2M+2}\,f_s^{(2M+1)}(N)
= \tfrac{B_{2M+2}}{2M+2}\,\bigl(-(s)_{2M+1}N^{-s-(2M+1)}\bigr)$ by \Cref{lem:iteratedDeriv_cpow}
(as $(-1)^{2M+1} = -1$). At $x = R$ the value is
$\tfrac{B_{2M+2}}{2M+2}\,f_s^{(2M+1)}(R)$, whose norm is
$\tfrac{|B_{2M+2}|}{2M+2}|(s)_{2M+1}|R^{-\operatorname{Re}s-(2M+1)} \to 0$ as $R \to \infty$,
since $\operatorname{Re} s + (2M+1) > 0$ (\texttt{Complex.norm\_cpow\_eq\_rpow\_re\_of\_pos},
\texttt{tendsto\_rpow\_neg\_atTop}). Hence in the limit
\[
  \int_N^{\infty} \widetilde{B}_{2M+1}\,f_s^{(2M+1)}
  = \frac{B_{2M+2}}{2M+2}\,(s)_{2M+1}\,N^{-s-(2M+1)}
    - \frac{1}{2M+2}\int_N^{\infty}\widetilde{B}_{2M+2}\,f_s^{(2M+2)} .
\]
Dividing by $(2M+1)!$ and using $(2M+2)\,(2M+1)! = (2M+2)!$ turns the first term into
$t_{M+1}(s,N)$; the second term is $-\frac{1}{(2M+2)!}\int_N^{\infty}\widetilde{B}_{2M+2}\,f_s^{(2M+2)}$,
which is exactly $R^{\infty}_{2M+2}(f_s;N)$ by \Cref{def:em_remainder_Ioi}. Hence
$R^{\infty}_{2M+1}(f_s;N) = t_{M+1}(s,N) + R^{\infty}_{2M+2}(f_s;N)$, the claim.
\end{proof}

\begin{lemma}[Holomorphy of the order-$p$ remainder integral]
\label{lem:rem_holo}
\lean{Complex.differentiableOn_eulerMaclaurinRemainderIoi_cpow}
\leanok
\uses{def:em_remainder_Ioi, lem:integrable_rem_zeta, lem:iteratedDeriv_cpow}
Fix integers $N \ge 1$ and $p \ge 1$. The map
$s \mapsto R^{\infty}_{p}(f_s; N)$ is holomorphic on the open half-plane
$\Omega_p := \{s : \operatorname{Re} s > 1 - p\}$.
\end{lemma}

\begin{proof}
\leanok
\uses{lem:integrable_rem_zeta, lem:iteratedDeriv_cpow}
By \Cref{lem:iteratedDeriv_cpow}, $f_s^{(p)}(x) = (-1)^p (s)_p\, x^{-s-p}$ on $x>0$, so
\[
  R^{\infty}_{p}(f_s; N)
  = \frac{(-1)^{p+1}}{p!}\,(-1)^p (s)_p \int_N^{\infty} \widetilde{B}_p(x)\, x^{-s-p}\,dx
  = -\frac{(s)_p}{p!}\,J_p(s), \qquad J_p(s) := \int_N^{\infty}\widetilde{B}_p(x)\,x^{-s-p}\,dx .
\]
The prefactor $-(s)_p/p! = -(\operatorname{ascPochhammer}\,p)(s)/p!$ is a polynomial in $s$,
hence entire; it remains to show $J_p$ is holomorphic on $\Omega_p$.

This is differentiation under the integral sign, identical in form to mathlib's holomorphy
of the Mellin transform (\texttt{mellin\_hasDerivAt\_of\_isBigO\_rpow}, which is itself an
application of \texttt{hasDerivAt\_integral\_of\_dominated\_loc\_of\_deriv\_le}). Fix
$s_0 \in \Omega_p$ and pick $\sigma_1$ with $1 - p < \sigma_1 < \operatorname{Re} s_0$; let
$V := \{s : \operatorname{Re} s > \sigma_1\}$, an open neighbourhood of $s_0$ contained in
$\Omega_p$, and choose a ball $\overline{B}(s_0, r) \subseteq V$. We verify the hypotheses
of \texttt{hasDerivAt\_integral\_of\_dominated\_loc\_of\_deriv\_le} on $\overline{B}(s_0,r)$
with integrand $G(s,x) = \widetilde{B}_p(x)\,x^{-s-p}$ and $s$-derivative
$\partial_s G(s,x) = -\widetilde{B}_p(x)\,x^{-s-p}\log x$:
\begin{itemize}
\item \emph{Integrability of $G(s,\cdot)$} for each $s$ in the ball: this is
  \Cref{lem:integrable_rem_zeta} (the kernel is bounded and
  $\|x^{-s-p}\| = x^{-\operatorname{Re}s - p}$ is integrable on $(N,\infty)$ because
  $\operatorname{Re} s + p > \sigma_1 + p > 1$); a.e.-strong-measurability of both $G$ and
  $\partial_s G$ in $x$ comes from \texttt{measurable\_fract} and continuity of
  $x \mapsto x^{-s-p}$ and $x \mapsto \log x$ on $(0,\infty)$.
\item \emph{$s$-derivative}: for $x > 0$ fixed, $s \mapsto x^{-s-p} = \exp(-(s+p)\log x)$ is
  entire with derivative $-x^{-s-p}\log x$ (\texttt{hasStrictDerivAt\_const\_cpow} for
  $z \mapsto x^{z}$, which has derivative $x^{z}\log x$, composed with the affine
  $s \mapsto -s-p$ via the chain rule), giving $\partial_s G(s,x)$ as stated.
\item \emph{Uniform domination}: with $C := \sup_{[0,1]}|B_p|$ the kernel bound
  (\Cref{lem:integrable_rem_zeta}), for all $s \in \overline{B}(s_0,r) \subseteq V$ and all
  $x \ge N$,
  \[
    \|\partial_s G(s,x)\| \le C\,|\log x|\,x^{-\operatorname{Re}s - p}
      \le C\,(\log x)\,x^{-\sigma_1 - p} =: h(x),
  \]
  and $h$ is integrable on $(N,\infty)$ since $\sigma_1 + p > 1$ (the logarithmic factor is
  absorbed: $\int_N^{\infty} x^{-a}\log x\,dx < \infty$ for $a > 1$; concretely dominate
  $\log x \le x^{\varepsilon}$ and use $\sigma_1 + p - \varepsilon > 1$ for small
  $\varepsilon$, \texttt{integrableOn\_Ioi\_rpow\_of\_lt}).
\end{itemize}
The lemma yields $\operatorname{HasDerivAt} J_p (\textstyle\int_N^\infty \partial_s G(s_0,x)\,dx)\,s_0$,
so $J_p$ is complex-differentiable at every $s_0 \in \Omega_p$; differentiability throughout
the open $\Omega_p$ gives holomorphy. Multiplying by the entire prefactor proves the claim.
\end{proof}

\begin{definition}[Analytic Euler--Maclaurin right-hand side for $\zeta$]
\label{def:zeta_em_rhs}
\lean{zetaEMExpansion}
\leanok
\uses{def:em_remainder_Ioi}
For integers $N \ge 1$, $M \ge 0$ and $s \in \mathbb{C}$, define
\[
  E_M(s,N) := \sum_{n=1}^{N-1} n^{-s} + \frac{N^{1-s}}{s-1} + \tfrac12 N^{-s}
    + \sum_{k=1}^{M} \frac{B_{2k}}{(2k)!}\,(s)_{2k-1}\,N^{-s-(2k-1)}
    + t_{M+1}(s,N) + R^{\infty}_{2M+2}(f_s; N),
\]
with $t_{M+1}(s,N) = \frac{B_{2M+2}}{(2M+2)!}(s)_{2M+1}N^{-s-(2M+1)}$ and
$R^{\infty}_{2M+2}(f_s;N)$ the order-$(2M+2)$ improper remainder
(\Cref{def:em_remainder_Ioi}). The last two summands together are the analytically continued
remainder $R_M(s,N) := t_{M+1}(s,N) + R^{\infty}_{2M+2}(f_s;N)$; by \Cref{lem:rem_ibp} they
equal $R^{\infty}_{2M+1}(f_s;N)$ wherever the latter converges
($\operatorname{Re} s > -2M$).
\end{definition}

\begin{lemma}[Analyticity of the remainder representation on $\operatorname{Re} s > -(2M+1)$]
\label{lem:zeta_em_analytic}
\lean{analyticOnNhd_zetaEMExpansion}
\leanok
\uses{def:zeta_em_rhs, lem:rem_holo, lem:iteratedDeriv_cpow}
For integer $N \ge 1$, the function $s \mapsto E_M(s,N)$ of \Cref{def:zeta_em_rhs} is
holomorphic on the punctured strip
$U_M := \{\operatorname{Re} s > -(2M+1)\}\setminus\{1\}$.
\end{lemma}

\begin{proof}
\leanok
\uses{def:zeta_em_rhs, lem:rem_holo}
$E_M(\cdot,N)$ is a finite sum of summands each holomorphic on $U_M$:
\begin{itemize}
\item each $n^{-s} = \exp(-s\log n)$ ($1 \le n \le N-1$) is entire, being a composition of
  the entire exponential with a linear map (\texttt{Complex.cpow} of a positive real base is
  entire in the exponent);
\item $N^{1-s}/(s-1)$ is holomorphic away from $s = 1$ (entire numerator over the nonzero
  denominator on $U_M$), and $\tfrac12 N^{-s}$ is entire;
\item each Bernoulli term $\frac{B_{2k}}{(2k)!}(s)_{2k-1}N^{-s-(2k-1)}$ ($1 \le k \le M$) and
  the term $t_{M+1}(s,N)$ are products of the polynomial $(s)_{j} = (\operatorname{ascPochhammer}\,j)(s)$
  with the entire exponential $N^{-s-j}$, hence entire;
\item the remainder $R^{\infty}_{2M+2}(f_s;N)$ is holomorphic on
  $\Omega_{2M+2} = \{\operatorname{Re} s > 1 - (2M+2)\} = \{\operatorname{Re} s > -(2M+1)\}
  \supseteq U_M$ by \Cref{lem:rem_holo} with $p = 2M+2$.
\end{itemize}
Sums and products of holomorphic functions are holomorphic, so $E_M(\cdot,N)$ is holomorphic
on $U_M$ (\texttt{DifferentiableOn.analyticOnNhd} packages this as analyticity on the open
set $U_M$).
\end{proof}

\begin{theorem}[Euler--Maclaurin expansion of $\zeta$]
\label{thm:zeta_em}
\lean{riemannZeta_eq}
\leanok
\uses{lem:zeta_em_conv, lem:zeta_em_analytic, lem:rem_ibp, def:zeta_em_rhs}
For $s \in \mathbb{C}$ with $\operatorname{Re} s > -(2M+1)$ and $s \neq 1$, and for
$N \ge 1$,
\[
  \zeta(s) = \sum_{n=1}^{N-1} n^{-s} + \frac{N^{1-s}}{s-1} + \tfrac12 N^{-s}
    + \sum_{k=1}^{M} \frac{B_{2k}}{(2k)!}\,(s)_{2k-1}\,N^{-s-(2k-1)} + R_{M}(s,N),
\]
equivalently $\zeta(s) = E_M(s,N)$, where $(s)_{j} = s(s+1)\cdots(s+j-1)$ is the rising
factorial, $f_s^{(j)}(x) = (-1)^j (s)_j\, x^{-s-j}$ (\Cref{lem:iteratedDeriv_cpow}), and
$R_M(s,N) = t_{M+1}(s,N) + R^{\infty}_{2M+2}(f_s;N)$ is the analytically continued remainder
(\Cref{def:zeta_em_rhs}).
\end{theorem}

\begin{proof}
\leanok
\uses{lem:zeta_em_conv, lem:zeta_em_analytic, lem:rem_ibp, def:zeta_em_rhs}
Work on the open set $U = \{\operatorname{Re} s > -(2M+1)\}\setminus\{1\}$. It is
preconnected: the half-plane $\{\operatorname{Re} s > -(2M+1)\}$ is convex
(\texttt{convex\_halfSpace\_re\_gt}), hence connected, and deleting the single point $s=1$
from a connected open subset of $\mathbb{C} \cong \mathbb{R}^2$ leaves it connected (a path
between two points may be perturbed off $1$). Both sides of the asserted identity are
holomorphic on $U$:
\begin{itemize}
\item $\zeta$ is holomorphic on $U$: it is differentiable at every $s \ne 1$
  (\texttt{differentiableAt\_riemannZeta});
\item $E_M(\cdot,N)$ is holomorphic on $U$ by \Cref{lem:zeta_em_analytic}.
\end{itemize}
On the sub-half-plane $W := \{\operatorname{Re} s > 1\} \subseteq U$ they agree: there
$\operatorname{Re} s > 1 > -2M$, so by \Cref{lem:rem_ibp} the continued remainder
$R_M(s,N) = t_{M+1}(s,N) + R^{\infty}_{2M+2}(f_s;N)$ equals the convergent order-$(2M+1)$
remainder $R^{\infty}_{2M+1}(f_s;N)$, whence $E_M(s,N)$ coincides with the right-hand side of
\Cref{lem:zeta_em_conv}, which equals $\zeta(s)$. Since $W$ is open and nonempty, $\zeta$ and
$E_M(\cdot,N)$ agree on a neighbourhood of any chosen $s_\star \in W$, i.e.\ they are
eventually equal at $s_\star$. The identity theorem
(\texttt{AnalyticOnNhd.eqOn\_of\_preconnected\_of\_eventuallyEq}, with the preconnected $U$,
both functions analytic on $U$, and eventual equality at $s_\star \in U$) propagates the
agreement to all of $U$. This is the claim for every $s \in U$, i.e.\ for
$\operatorname{Re} s > -(2M+1)$, $s \ne 1$.
\end{proof}

\section{The Olver--Borwein error bound}
\label{sec:em_error_bound}

This section proves the sharp tail estimate \Cref{thm:zeta_error_bound}, the bound
implemented as \texttt{zeta\_error\_bound} in \texttt{python/src/euler\_maclaurin/zeta\_values.py}. Throughout write
$m := M+1$, so the relevant even index is $2m = 2M+2$ and the convergence strip
$\operatorname{Re} s > -(2M+1)$ reads $\operatorname{Re} s + 2m > 1$. The remainder of
\Cref{thm:zeta_em} is, by \Cref{def:zeta_em_rhs},
\[
  R_M(s,N) \;=\; t_{M+1}(s,N) + R^{\infty}_{2M+2}(f_s; N),
  \qquad
  t_{M+1}(s,N) = \frac{B_{2M+2}}{(2M+2)!}\,(s)_{2M+1}\,N^{-s-(2M+1)} .
\]

The triangle inequality on this split is insufficient: bounding
$|R^{\infty}_{2M+2}| \le K\,|t_{M+1}|$ only gives $|R_M| \le (1+K)|t_{M+1}|$, weaker than the
claim. The sharp constant comes from first collapsing the split back into a \emph{single}
integral of the \emph{centered} kernel $\widetilde{B}_{2m} - B_{2m}$, whose period-mean is
exactly $-B_{2m}$, and then exploiting that this centered kernel has constant sign with the
oscillatory part beating the weight $x^{-a}$ in the right direction.

\begin{lemma}[Centered-kernel representation of the remainder]
\label{lem:rem_centered}
\lean{zetaRemainder_eq_integral_bernoulliFun_fract_sub}
\leanok
\uses{def:zeta_em_rhs, def:em_remainder_Ioi, lem:iteratedDeriv_cpow, lem:integrable_rem_zeta, def:periodizedBernoulli, def:bernoulli}
Let $N \ge 1$, $M \ge 0$, $s \in \mathbb{C}$ with $\operatorname{Re} s > -(2M+1)$, and set
$m := M+1$. Then
\[
  R_M(s,N) \;=\; -\frac{(s)_{2m}}{(2m)!}\int_N^{\infty}
      \bigl(\widetilde{B}_{2m}(x) - B_{2m}\bigr)\, x^{-s-2m}\,dx,
\]
where $\widetilde{B}_{2m}(x) = \operatorname{bernoulliFun}(2m)\,(\operatorname{Int.fract} x)$ and
the integral converges absolutely (\Cref{lem:integrable_rem_zeta}, since
$\operatorname{Re} s + 2m > 1$).
\end{lemma}

\begin{proof}
\leanok
\uses{def:zeta_em_rhs, def:em_remainder_Ioi, lem:iteratedDeriv_cpow, lem:integrable_rem_zeta}
By \Cref{lem:iteratedDeriv_cpow}, $f_s^{(2m)}(x) = (-1)^{2m}(s)_{2m}\,x^{-s-2m}
= (s)_{2m}\,x^{-s-2m}$, so by \Cref{def:em_remainder_Ioi} (with $p = 2m$, $(-1)^{2m+1} = -1$)
\[
  R^{\infty}_{2M+2}(f_s; N) = -\frac{1}{(2m)!}\int_N^{\infty}\widetilde{B}_{2m}(x)\,f_s^{(2m)}(x)\,dx
  = -\frac{(s)_{2m}}{(2m)!}\int_N^{\infty}\widetilde{B}_{2m}(x)\,x^{-s-2m}\,dx,
\]
the integral converging absolutely for $\operatorname{Re} s + 2m > 1$
(\Cref{lem:integrable_rem_zeta}). On that strip both $\int_N^{\infty}\widetilde{B}_{2m}(x)x^{-s-2m}dx$
and $\int_N^{\infty} x^{-s-2m}dx$ converge, so we may split
$\widetilde{B}_{2m} = B_{2m} + (\widetilde{B}_{2m} - B_{2m})$ linearly. The constant part is
evaluated by the fundamental theorem of calculus,
\[
  \int_N^{\infty} x^{-s-2m}\,dx = \frac{N^{1-s-2m}}{s+2m-1}
  \qquad(\text{antiderivative } \tfrac{x^{1-s-2m}}{1-s-2m};\ \text{boundary at } \infty
  \text{ vanishes as } \operatorname{Re}(1-s-2m) < 0),
\]
and using $\dfrac{(s)_{2m}}{s+2m-1} = (s)_{2m-1} = (s)_{2M+1}$ together with $1-2m = -(2M+1)$,
\[
  -\frac{(s)_{2m}}{(2m)!}\,B_{2m}\,\frac{N^{1-s-2m}}{s+2m-1}
  = -\frac{B_{2M+2}}{(2M+2)!}\,(s)_{2M+1}\,N^{-s-(2M+1)} = -t_{M+1}(s,N).
\]
Hence $R^{\infty}_{2M+2}(f_s;N) = -t_{M+1}(s,N) - \frac{(s)_{2m}}{(2m)!}
\int_N^{\infty}(\widetilde{B}_{2m}-B_{2m})\,x^{-s-2m}\,dx$. Adding $t_{M+1}(s,N)$, the two
$t_{M+1}$ terms cancel and the claim follows from $R_M = t_{M+1} + R^{\infty}_{2M+2}$
(\Cref{def:zeta_em_rhs}).
\end{proof}

\begin{lemma}[Sign and cosine--Fourier expansion of the centered Bernoulli kernel]
\label{lem:bernoulli_centered_fourier}
\lean{bernoulliFun_sub_self_eq_tsum_cos}
\leanok
\uses{def:bernoulliFun, def:periodizedBernoulli, def:bernoulli}
Let $m \ge 1$ and $A_m := \dfrac{2\,(2m)!}{(2\pi)^{2m}} > 0$. For all $x \in \mathbb{R}$,
\[
  \widetilde{B}_{2m}(x) - B_{2m}
    = (-1)^{m+1} A_m \sum_{k=1}^{\infty}\frac{\cos(2\pi k x) - 1}{k^{2m}} .
\]
Consequently the centered kernel has constant sign $(-1)^m$ (i.e.
$(-1)^m\bigl(\widetilde{B}_{2m}(x) - B_{2m}\bigr) \ge 0$ for all $x$), and
\[
  \bigl|\widetilde{B}_{2m}(x) - B_{2m}\bigr|
    = A_m \sum_{k=1}^{\infty}\frac{1 - \cos(2\pi k x)}{k^{2m}} ,
\]
each summand being nonnegative.
\end{lemma}

\begin{proof}
\leanok
\uses{def:bernoulliFun, def:periodizedBernoulli, def:bernoulli}
For $x \in [0,1]$ mathlib's Fourier expansion of the Bernoulli polynomials
(\texttt{hasSum\_one\_div\_nat\_pow\_mul\_cos}, the cosine specialisation of
\texttt{hasSum\_one\_div\_nat\_pow\_mul\_fourier}, valid as $2m \ge 2$) gives
\[
  \sum_{n=1}^{\infty}\frac{1}{n^{2m}}\bigl(e^{2\pi i n x} + e^{-2\pi i n x}\bigr)
  = -\frac{(2\pi i)^{2m}}{(2m)!}\,B_{2m}(x).
\]
The bracket is $2\cos(2\pi n x)$ and $(2\pi i)^{2m} = (-1)^m (2\pi)^{2m}$, so
\[
  B_{2m}(x) = (-1)^{m+1}\,\frac{2\,(2m)!}{(2\pi)^{2m}}\sum_{n=1}^{\infty}\frac{\cos(2\pi n x)}{n^{2m}}
            = (-1)^{m+1} A_m \sum_{n=1}^{\infty}\frac{\cos(2\pi n x)}{n^{2m}}, \qquad x \in [0,1].
\]
Taking $x = 0$ (where $\cos = 1$ and $\sum n^{-2m} = \zeta(2m)$) gives
$B_{2m} = B_{2m}(0) = (-1)^{m+1} A_m\,\zeta(2m)$. Subtracting and using
$\widetilde{B}_{2m}(x) = B_{2m}(\operatorname{Int.fract} x)$ with $\operatorname{Int.fract} x \in [0,1)$
(and periodicity, both sides being $1$-periodic) gives the displayed identity for all
$x \in \mathbb{R}$. Since $\cos(2\pi k x) - 1 \le 0$ for every $k$, the series has sign
$(-1)^{m+1}\cdot(\le 0) = (-1)^m\cdot(\ge 0)$, whence the constant-sign claim; and
$|\widetilde{B}_{2m}(x) - B_{2m}| = A_m\sum_k (1-\cos(2\pi k x))/k^{2m}$ because each term
$1 - \cos(2\pi k x) \ge 0$.
\end{proof}

\begin{lemma}[Period-mean of the centered kernel]
\label{lem:bernoulli_centered_mean}
\lean{integral_abs_bernoulliFun_sub_self}
\leanok
\uses{lem:bernoulli_centered_fourier, def:bernoulliFun}
For $m \ge 1$, $\displaystyle\int_0^1 \bigl|\widetilde{B}_{2m}(x) - B_{2m}\bigr|\,dx = |B_{2m}|$.
(In Lean the integrand uses the Bernoulli polynomial $B_{2m}(x)$, which coincides with
$\widetilde{B}_{2m}(x)$ on $[0,1)$, hence the integrals agree.)
\end{lemma}

\begin{proof}
\leanok
\uses{lem:bernoulli_centered_fourier, def:bernoulliFun}
The centered kernel has constant sign (\Cref{lem:bernoulli_centered_fourier}), so
\[
  \int_0^1 \bigl|\widetilde{B}_{2m}(x) - B_{2m}\bigr|\,dx
  = \left|\int_0^1 \bigl(\widetilde{B}_{2m}(x) - B_{2m}\bigr)\,dx\right|
  = \left|\int_0^1 B_{2m}(x)\,dx - B_{2m}\right| = |0 - B_{2m}| = |B_{2m}|,
\]
using $\widetilde{B}_{2m} = B_{2m}\circ\operatorname{Int.fract}$ agreeing with $B_{2m}$ on $[0,1)$
and the vanishing mean $\int_0^1 B_{2m} = 0$ (\texttt{integral\_bernoulliFun\_eq\_zero}, as
$2m \ne 0$).
\end{proof}

\begin{lemma}[Pólya positivity of the cosine integral]
\label{lem:polya_cos}
\lean{integral_Ioi_cos_mul_rpow_nonneg}
\leanok
For an integer $N \ge 1$, a real $a > 1$, and an integer $k \ge 1$,
\[
  \int_N^{\infty}\cos(2\pi k x)\,x^{-a}\,dx \;\ge\; 0 .
\]
\end{lemma}

\begin{proof}
\leanok
This is the one new analytic input (Pólya, 1918): the cosine transform of a
nonnegative decreasing function is nonnegative.

In Lean the argument is run on the conjugate \emph{sine} integral, which is more convenient
because $\sin(2\pi k x)$ \emph{vanishes} at $x=N$ (as $kN \in \mathbb{Z}$), so its sign-constant
half-periods $[N+\tfrac{j}{2k}, N+\tfrac{j+1}{2k}]$ start exactly at $N$. Integration by parts on
$[N,R]$ gives
\[
  \int_N^{R}\cos(2\pi k x)\,x^{-a}\,dx
  = \frac{\sin(2\pi k R)}{2\pi k}R^{-a}
    + \frac{a}{2\pi k}\int_N^{R}\sin(2\pi k x)\,x^{-(a+1)}\,dx,
\]
the boundary term at $N$ vanishing since $\sin(2\pi k N)=0$; letting $R\to\infty$ (the boundary
term at $R$ is $O(R^{-a})\to 0$) yields
$\int_N^{\infty}\cos(2\pi k x)x^{-a}dx = \frac{a}{2\pi k}\int_N^{\infty}\sin(2\pi k x)x^{-(a+1)}dx$,
so it suffices to prove $\int_N^{\infty}\sin(2\pi k x)\,x^{-(a+1)}\,dx \ge 0$ (formalized as
\texttt{integral\_Ioi\_sin\_mul\_rpow\_nonneg}). Write $\psi(x) := x^{-(a+1)}$, positive,
decreasing, and integrable on $(N,\infty)$ (as $a+1 > 1$), so $\psi\sin(2\pi k\cdot)$ is absolutely
integrable. Split $(N,\infty)$ into the half-periods
$I_j := \bigl[N + \tfrac{j}{2k}, N + \tfrac{j+1}{2k}\bigr]$, $j \ge 0$; since $kN \in \mathbb{Z}$,
on $I_j$ we have $\sin(2\pi k x) = (-1)^j\bigl|\sin(2\pi k x)\bigr|$ (constant sign $(-1)^j$,
starting positive on $I_0$ because $\sin$ vanishes at the left endpoint $N$). Put
\[
  c_j := \int_{I_j}\psi(x)\,|\sin(2\pi k x)|\,dx \ge 0,
  \qquad\text{so}\qquad
  \int_{N+\frac{j}{2k}}^{N+\frac{j+1}{2k}}\sin(2\pi k x)\,\psi(x)\,dx = (-1)^j c_j .
\]
The shift $x \mapsto x + \tfrac{1}{2k}$ maps $I_j$ to $I_{j+1}$ and preserves $|\sin(2\pi k x)|$
(as $\sin(\theta+\pi) = -\sin\theta$), so
$c_{j+1} = \int_{I_j}\psi\bigl(x+\tfrac1{2k}\bigr)|\sin(2\pi k x)|\,dx \le c_j$ because $\psi$ is
decreasing. Hence the truncated integrals $\int_N^{N+\frac{n}{2k}}\sin(2\pi k x)\psi(x)\,dx
= \sum_{j<n}(-1)^j c_j$ are the partial sums of an alternating series with antitone nonnegative
terms; they converge to the improper integral (by
\texttt{intervalIntegral\_tendsto\_integral\_Ioi}, the upper limit $N+\tfrac{n}{2k}\to\infty$).
By \texttt{Antitone.alternating\_series\_le\_tendsto} the empty even partial sum ($n=0$) is
$\le$ the limit, i.e. $0 \le \int_N^{\infty}\sin(2\pi k x)\psi(x)\,dx$.
\end{proof}

\begin{lemma}[Weighted averaging bound for the centered kernel]
\label{lem:centered_kernel_avg}
\lean{integral_Ioi_abs_bernoulliFun_fract_sub_mul_rpow_le}
\leanok
\uses{lem:bernoulli_centered_fourier, lem:bernoulli_centered_mean, lem:polya_cos}
For $m \ge 1$, integer $N \ge 1$, and real $a > 1$,
\[
  \int_N^{\infty}\bigl|\widetilde{B}_{2m}(x) - B_{2m}\bigr|\,x^{-a}\,dx
    \;\le\; |B_{2m}| \int_N^{\infty} x^{-a}\,dx .
\]
\end{lemma}

\begin{proof}
\leanok
\uses{lem:bernoulli_centered_fourier, lem:bernoulli_centered_mean, lem:polya_cos}
By \Cref{lem:bernoulli_centered_fourier}, $|\widetilde{B}_{2m}(x) - B_{2m}|
= A_m \sum_{k\ge1}(1-\cos(2\pi k x))/k^{2m}$ with all summands nonnegative, where
$A_m\,\zeta(2m) = |B_{2m}|$ (\Cref{lem:bernoulli_centered_mean}, equivalently from the $x=0$
value in \Cref{lem:bernoulli_centered_fourier}). Multiplying by the nonnegative weight
$x^{-a}$ and integrating termwise on $(N,\infty)$ (justified by Tonelli, all terms being
nonnegative and measurable),
\[
  \int_N^{\infty}\bigl|\widetilde{B}_{2m} - B_{2m}\bigr|\,x^{-a}\,dx
  = A_m \sum_{k=1}^{\infty}\frac{1}{k^{2m}}\int_N^{\infty}\bigl(1 - \cos(2\pi k x)\bigr)x^{-a}\,dx .
\]
Splitting the inner integral and using $\sum_k k^{-2m} = \zeta(2m)$,
\[
  = A_m\,\zeta(2m)\int_N^{\infty} x^{-a}\,dx
    - A_m \sum_{k=1}^{\infty}\frac{1}{k^{2m}}\int_N^{\infty}\cos(2\pi k x)\,x^{-a}\,dx
  \;\le\; |B_{2m}|\int_N^{\infty} x^{-a}\,dx,
\]
since each $\int_N^{\infty}\cos(2\pi k x)\,x^{-a}\,dx \ge 0$ by \Cref{lem:polya_cos}, so the
subtracted sum is $\ge 0$. (Both $\int_N^{\infty} x^{-a}dx = N^{1-a}/(a-1)$ and the rearranged
series converge, $a > 1$.)
\end{proof}

\begin{theorem}[Borwein error bound]
\label{thm:zeta_error_bound}
\lean{norm_zetaRemainder_le}
\leanok
\uses{thm:zeta_em, lem:rem_centered, lem:centered_kernel_avg, lem:iteratedDeriv_cpow}
With the notation of \Cref{thm:zeta_em} and $\operatorname{Re} s > -(2M+1)$, the
remainder satisfies
\[
  \bigl| R_M(s, N) \bigr| \;\le\;
    \left| \frac{s + 2M+1}{\operatorname{Re} s + 2M + 1} \right| \cdot \bigl| t_{M+1} \bigr|,
  \qquad
  t_{M+1} = \frac{B_{2M+2}}{(2M+2)!}\,(s)_{2M+1}\,N^{-s-(2M+1)},
\]
i.e. the tail is controlled by the first omitted Bernoulli term scaled by the explicit
constant $K := \bigl|\tfrac{s+2M+1}{\operatorname{Re} s+2M+1}\bigr| \ge 1$ (note
$K = |(s)_{2M+2}| / (|(s)_{2M+1}|\,(\operatorname{Re} s+2M+1))$, since
$(s)_{2M+2} = (s)_{2M+1}(s+2M+1)$). This is the bound implemented as \texttt{zeta\_error\_bound} in
\texttt{python/src/euler\_maclaurin/zeta\_values.py}.
\end{theorem}

\begin{proof}
\leanok
\uses{lem:rem_centered, lem:centered_kernel_avg, lem:iteratedDeriv_cpow}
Write $m := M+1$, so $2m = 2M+2$ and $a := \operatorname{Re} s + 2m = \operatorname{Re} s + 2M+2
> 1$ (as $\operatorname{Re} s > -(2M+1)$). By the centered-kernel representation
(\Cref{lem:rem_centered}) and $\|x^{-s-2m}\| = x^{-\operatorname{Re} s - 2m} = x^{-a}$ for
$x > 0$ (\texttt{Complex.norm\_cpow\_eq\_rpow\_re\_of\_pos}), the triangle inequality for the
Bochner integral gives
\[
  \bigl|R_M(s,N)\bigr|
  = \frac{|(s)_{2m}|}{(2m)!}\left|\int_N^{\infty}\bigl(\widetilde{B}_{2m}(x)-B_{2m}\bigr)x^{-s-2m}\,dx\right|
  \le \frac{|(s)_{2m}|}{(2m)!}\int_N^{\infty}\bigl|\widetilde{B}_{2m}(x)-B_{2m}\bigr|\,x^{-a}\,dx .
\]
(Here the complex oscillation of $x^{-s-2m}$ plays no role: it is absorbed by the modulus,
leaving the real exponent $a$.) Applying the weighted averaging bound
(\Cref{lem:centered_kernel_avg}) and $\int_N^{\infty} x^{-a}dx = N^{1-a}/(a-1)$,
\[
  \bigl|R_M(s,N)\bigr|
  \le \frac{|(s)_{2m}|}{(2m)!}\,|B_{2m}|\,\frac{N^{1-a}}{a-1}
  = \frac{|(s)_{2M+2}|\,|B_{2M+2}|}{(2M+2)!}\cdot\frac{N^{-\operatorname{Re} s-(2M+1)}}{\operatorname{Re} s+2M+1},
\]
using $a - 1 = \operatorname{Re} s + 2M+1$ and $1 - a = -\operatorname{Re} s - (2M+1)$. Finally,
expanding $K\,|t_{M+1}|$ with $|(s)_{2M+2}| = |(s)_{2M+1}|\,|s+2M+1|$,
\[
  K\,|t_{M+1}|
  = \frac{|s+2M+1|}{\operatorname{Re} s+2M+1}\cdot\frac{|B_{2M+2}|}{(2M+2)!}\,|(s)_{2M+1}|\,N^{-\operatorname{Re} s-(2M+1)}
  = \frac{|(s)_{2M+2}|\,|B_{2M+2}|}{(2M+2)!}\cdot\frac{N^{-\operatorname{Re} s-(2M+1)}}{\operatorname{Re} s+2M+1},
\]
which is exactly the right-hand side above. Hence $|R_M(s,N)| \le K\,|t_{M+1}|$, as claimed.
\end{proof}

\begin{remark}[log-Gamma and Hardy's $Z$]
\label{rem:log_gamma}
The companion results for $\log\Gamma$ (Euler--Maclaurin / Stirling expansion) and
for Hardy's $Z(t) = e^{i\theta(t)}\zeta(\tfrac12+it)$ with the Riemann--Siegel
$\theta$ follow the same template: apply \Cref{thm:euler_maclaurin} to
$f(x) = \log x$ and track the analogous remainder bound. These are implemented in
\texttt{python/src/euler\_maclaurin/gamma\_values.py} and
\texttt{python/src/euler\_maclaurin/zeta\_zeros.py}.
\end{remark}
